\section{Introduction}

The Internet of Things (IoT) has become an integral component of contemporary society, with Consumer IoT (CIoT) devices accounting for nearly 50\% of the overall IoT market \cite{cisco2020cisco}. Although these technologies provide substantial benefits in everyday contexts, their pervasive deployment across domestic, professional, and public environments—together with their persistent connection to network infrastructures—raises significant concerns related to privacy and security. In this context, smart home devices are of particular interest, as they continuously and often unobtrusively collect data related to users’ daily activities. Such characteristics make them a valuable source for IoT forensics, enabling the acquisition and analysis of device-generated data, including network traffic, which can be leveraged to infer events, behaviors, or usage patterns \cite{stoyanova}. Two communication protocols are predominantly used for interconnecting smart home devices: Wi-Fi and Zigbee. While a substantial body of literature has demonstrated the feasibility of IoT forensic tools for capturing and analyzing traffic generated by Wi-Fi–enabled devices \cite{fabio,fabio2,fabio3}, forensic methodologies specifically targeting Zigbee-based devices remain comparatively underexplored. This gap is largely due to the inherent challenges associated with Zigbee traffic analysis. First, the characteristics of the Zigbee protocol, which departs significantly from the conventional TCP/IP stack, introduce additional complexity in the interpretation of captured traffic. Furthermore, Zigbee networks typically generate smaller packet sizes and lower traffic volumes than Wi-Fi, limiting the amount of observable information available for forensic inference. Finally, Zigbee commonly relies on a mesh network topology, in which devices forward traffic on behalf of others rather than communicating directly with a central access point, as is the case with Wi-Fi. This architectural feature further complicates traffic analysis, as relay operations introduce noise into traffic traces and such noise is inherently time-varying due to dynamic routing path adaptations. As a result, forensic insights derived from traffic analysis in a given Zigbee network topology may not generalize to other topologies, in contrast to what is more commonly observed in Wi-Fi–based environments.
This study addresses the aforementioned challenges through the following main contributions:
\begin{itemize}
    \item \textit{ZIOTP2026 traffic dataset\footnote{https://github.com/antonio-boiano/ZIOTP2026}:} 
    We collected and publicly released a comprehensive dataset of Zigbee network traffic captured from 21 off-the-shelf smart home devices deployed under two distinct network configurations. The dataset includes both background traffic and traffic generated through controlled user interactions, and is fully labeled to support IoT forensic tasks such as device identification, event detection, and traffic classification. To the best of the authors’ knowledge, this represents the largest publicly available labeled Zigbee traffic dataset specifically designed for IoT forensics.

    \item \textit{Traffic analysis:} 
    Leveraging a \gls*{soa}-based methodology, we developed a \gls*{ml} framework for Zigbee device classification, addressing two core tasks: device type classification and individual device identification. To mitigate the complexity introduced by multi-hop communications in mesh network topologies—where devices forward packets on behalf of others—we restrict the analysis to first-hop messages, thereby reducing bias due to retransmission and relaying behavior. Nevertheless, we show that temporal traffic features remain strongly affected by topology-dependent dynamics. Experimental results demonstrate that, while intra-topology classification achieves excellent performance (with F1-scores above 0.95), cross-topology generalization suffers from significant degradation. This highlights the critical and often overlooked impact of Zigbee mesh topology on the generalizability of traffic-based classification models.

    \item \textit{Storage--accuracy analysis:} 
    We investigate the trade-off between traffic storage requirements and the accuracy of downstream forensic analysis tasks. Specifically, we compare lossless traffic compression—obtained by applying state-of-the-art compression algorithms directly to raw Zigbee traffic—with lossy compression of traffic representations, where scalar quantization is applied to the statistical features used as input to \gls*{ml} models. Our results show that the latter approach enables a substantial reduction in storage bitrate without degrading classification accuracy, making it a promising solution for IoT forensic systems tasked with long-term storage of large-scale traffic traces.
\end{itemize}

% \textcolor{red}{
% The remainder of this paper is organized as follows:
% In section \ref{background}, the background for the ZigBee protocols is provided,
% Section \ref{state-of-the-art} presents a discussion of the main related works in IoT forensics and ZigBee network traffic analysis, Section \ref{ZIOTP2026} describes the ZIOTP2026 dataset and its creation methodology, Section \ref{methodology} details the feature extraction and data preprocessing techniques adopted, Sections \ref{iat-cls} and \ref{feat-cls} showcase the performance of state-of-the-art methodologies in multi \gls*{pan} domain. Section \ref{deepl} introduces a novel model for traffic analysis and present its performance across multiple \gls*{pan} domain, and lastly Section \ref{conclusions} conclude the work with final remarks and directions for future research.
% }

The remainder of this paper is organized as follows. 
Section \ref{background} introduces the Zigbee protocol architecture, while Section \ref{state-of-the-art} reviews related work in IoT forensics and Zigbee traffic analysis. 
Section \ref{ZIOTP2026} presents the ZIOTP2026 dataset and the adopted data collection methodology. 
Section \ref{sec:methodology} describes the feature extraction and classification framework, and Section \ref{results} discusses the experimental results under both intra-topology and cross-topology settings. 
Section \ref{sec:compression} analyzes the trade-off between traffic storage and classification accuracy. 
Finally, Section \ref{sec:conclusion} concludes the paper and outlines directions for future research.